\documentclass[11pt]{article}
\usepackage[a4paper,margin=1in]{geometry}
\usepackage[T1]{fontenc}
\usepackage[utf8]{inputenc}
\usepackage{lmodern}
\usepackage[english]{babel}
\usepackage{microtype}
\usepackage{amsmath,amssymb,mathtools}
\usepackage{enumitem}
\setlist[enumerate]{leftmargin=2em,itemsep=0.25em,topsep=0.25em}
\setlength{\parindent}{0pt}
\setlength{\parskip}{0.55em}
\newcommand{\DR}{\mathrm{DR}}
\newcommand{\Gr}{\operatorname{Gr}}
\newcommand{\RG}{R\Gamma}
\newcommand{\HH}{\mathcal H}
\newcommand{\C}{\mathbb C}
\begin{document}

\begin{flushright}
28 October 1976
\end{flushright}

Dear Luc,

Here is a global proof of the existence of a \(\DR\) complex giving rise to mixed Hodge theory.

\paragraph{Objects considered.}
The derived category of the following category of bifiltered complexes \((K^\ast,d,W,F)\):
\[
\begin{array}{rcl}
K^i &:& \text{a quasi-coherent sheaf, as are } W_jK^i \text{ and } F^jK^i,\\
d &:& \text{a first-order differential operator, respecting } W \text{ and } F,\\
\Gr_F(d) &:& \mathcal O\text{-linear.}
\end{array}
\]

\paragraph{Variants.}
\begin{enumerate}[label=\alph*)]
\item Require only that \(\HH^i\Gr_FK\) be quasi-coherent.
\item In this letter we restrict ourselves to \(F\), rather than to \(W\).
\end{enumerate}

\paragraph{Remarks.}
\begin{enumerate}[label=\alph*)]
\item There are functors \(Rf_\ast\).

\item For \(X_\bullet \xrightarrow{\varepsilon} X\), define \(R\varepsilon_\ast\) as follows. After resolving and applying \(\varepsilon_\ast\), one obtains a double complex \(K^{ij}\) on \(X\), with \(W,F\), where \(i\) is the index of \(X_i\) and \(j\) is the degree. One takes the associated simple complex; for \(F\), the simple complex associated to the \(F^i\), and, for \(W\), the diagonal between \(i\) and \(W\):
\[
  W^i=\bigoplus_{i'+i''=i} W^{i'}K^{i'',\ast}.
\]

\item There is inverse image by a smooth morphism \(f:Y\to X\):
\[
  f^\ast K\; {:=}_{\mathrm{dfn}}\; \Omega_Y^\ast\otimes_{f^{-1}\Omega_X^\ast} f^{-1}K.
\]
\end{enumerate}

\paragraph{What one wants.}
Construct functorially, for \(U\subset X\), with \(X\) proper, such an object on \(X\), whose \(R\Gamma\) ``is'' a mixed Hodge complex and gives rise to the mixed Hodge structure on \(H(U,\C)\). In particular, the assertions of 8.19 in Hodge III are true.

I realize that asking for this is too much: only \((K^\ast,d,F)\) can be canonical in this way. Why?

The construction is given in Hodge III, 8.3: take a resolution
\[
  (U_\bullet,X_\bullet) \xrightarrow{\varepsilon} (U,X)
\]
as in Hodge III, and, for \(D_\bullet=X_\bullet-U_\bullet\),
\[
  R\varepsilon_\ast\bigl(\Omega^\ast_{X_\bullet}(\log D_\bullet),W,F\bigr).
\]

The problem is the independence of the choice of \(U_\bullet,X_\bullet\).

Thus let
\[
\begin{array}{ccc}
(U'_\bullet,X'_\bullet) & \xrightarrow{\varphi} & (U_\bullet,X_\bullet)\\
\big\downarrow{\varepsilon'} && \big\downarrow{\varepsilon}\\
(U,X) && (U,X)
\end{array}
\]
and
\[
  K=R\varepsilon_\ast\Omega(\log),\qquad
  K'=R\varepsilon'_\ast\Omega(\log),\qquad
  \varphi^\ast:K\longrightarrow K'
\]
on \(X\), proper. We have:
\begin{enumerate}[label=\arabic*)]
\item \(\RG(X,K)\xrightarrow{\sim}\RG(X,K')\). This is \(H^\ast(U,\C)\).
\item \(\RG\bigl(X,\Gr_F^iK\bigr)\xrightarrow{\sim}\RG\bigl(X,\Gr_F^iK'\bigr)\). This is \(\Gr_F^iH^\ast(U,\C)\), since, for \(F\), \(E_1=E_\infty\).
\end{enumerate}

Now one uses the method of SGA \(4\frac12\), finiteness. Proceeding by induction on \(\dim X\), and noting that what is true over \(\C\) is also true over any field of characteristic \(0\), one finds that
\[
  \Gr_F^i(K)\longrightarrow \Gr_F^i(K')
\]
is a quasi-isomorphism away from finitely many points, rather away from a skyscraper, and that the third vertex of the triangle has zero cohomology. Restricting to \(U\), it is a quasi-isomorphism everywhere.

\paragraph{Conclusion.}
For \(U/\C\), there is a complex \((K^\ast,d,F)\), functorial up to quasi-isomorphism, with
\[
\begin{cases}
\HH\bigl(K^{\ast,\mathrm{an}}\bigr)=\C_U,
& \mathbb H(U,K^\ast)=\mathbb H(U^{\mathrm{an}},K^{\ast,\mathrm{an}})=H(U,\C),\\[0.35em]
\text{for }U\text{ proper, the spectral sequence of }F\text{ degenerates at }E_1,
& \text{and abuts to the Hodge filtration.}
\end{cases}
\]

\begin{flushright}
Yours,\\
P. Deligne
\end{flushright}

\end{document}
